\uuid{9j9h}
\titre{Congruences}
\auteur{bodin ; croizat ; sacre}
\organisation{exo7}
\licence{CC BY-NC-SA 4.0}
\datecreate{2022-03-08}
\niveau{L1}
\module{Algèbre}
\chapitre{Arithmétique dans Z}
\sousChapitre{Autre}
\difficulte{1}
\theme{Congruences}
\exo7qcmid{hilisit.arithmetique.041}

\contenu{
\begin{question}
\qclassification{103.99}
    Quelles sont les affirmations vraies ?
    \begin{answers}
    \bad{$456\,789 \equiv 0 \; [2]$}
    \good{$43\,210 \equiv 0 \; [5]$}
    \bad{$23\,769 \equiv 3 \; [9]$}
    \bad{$10\,326 \equiv 8 \; [10]$}
    \end{answers}
    \begin{explanations}
    $456\,789$ est impair, donc n’est pas congru à $0$ modulo $2$. \\
    $43\,210$ est divisible par $5$ donc congru à $0$ modulo $5$. \\
    $23\,769$ est divisible par $9$ (la somme des chiffres est divisible par $9$) donc congru à $0$ modulo $9$. \\
    $10\,326 \equiv 6 \; [10]$, réduire modulo $10$ c'est garder le chiffre des unités.
    \end{explanations}
\end{question}
}
